(1) 周期相位因子的傅里叶展开\(\exp(im\cos\theta)=\sum_{n=-\infty}^\infty i^nJ_n(m)e^{in\theta}\)。因此透射波可分解为横向波数\(k_{x,n}=2\pi n/d\)的平面波，每级复幅\(a_n=i^nJ_n(m)\)；该结果来自振幅的相干傅里叶分解，而不是对入射强度直接调制。
(2) 若入射、出射介质折射率均为1，则级次\(n\)满足\(\sin\theta_n=n\lambda/d\)，仅\(|n\lambda/d|\le1\)的级次为远场传播模；在所有衍射级均可传播、并适当使用功率归一化的理想近轴模型中，级次功率分数为\(P_n/P_{\rm in}=J_n(m)^2\)，且\(\sum_{n=-\infty}^\infty J_n(m)^2=1\)。若部分级次倏逝，不能把所有\(J_n^2\)无修正地解释为远场输送功率。
(3) 当\(m\ll1\)，\(J_0(m)=1-m^2/4+O(m^4)\)、\(J_1(m)=m/2+O(m^3)\)，故零级强度\(J_0^2=1-m^2/2+O(m^4)\)，一级各自\(J_{\pm1}^2=m^2/4+O(m^4)\)；三项合计\(1+O(m^4)\)。在二阶弱调制下，零级损失的能量恰转入对称的一级衍射，不凭空产生能量。
(4) 要消除零级，须\(J_0(m)=0\)，其首个正根\(m_1\simeq2.404826\)。此时大约\(J_{\pm1}(m_1)^2\simeq0.2695\)分别进入\(\pm1\)级，两者总和仅约0.539，余下能量进入\(\pm2,\pm3,\ldots\)高阶；因此零级完全抑制不意味着一级级次总效率为100%。有限空间口径和光栅色散还会改变实际角谱。
